\documentclass[12pt]{article}

\usepackage[portuguese]{babel}
\usepackage{float}
\usepackage{listings}
\usepackage{indentfirst}
\usepackage{amsfonts, amssymb, amsmath}
\usepackage{graphicx}
\usepackage{tikz}
\usepackage{enumitem}
\usepackage[colorlinks=true, allcolors=blue]{hyperref}
\usepackage[letterpaper,
            top=1.5cm,
            bottom=1.5cm,
            left=2.75cm,
            right=2.75cm,
            marginparwidth=1.75cm]{geometry}
        

\begin{document}
\begin{titlepage}
    \centering
    \vspace*{1cm}
    
    % 1. Institution/Department
    \Large\textbf{Universidade Estadual de Maringá} \\
    \large Ciência da Computação \\
    \vspace{5cm}
    
    % 3. Main Title
    \Huge\textbf{Relatório do Trabalho de Programação de Sistemas} \\
    \vspace{5cm}
    
    % 4. Author and Supervisor Info
    \begin{minipage}{0.4\textwidth}
        \begin{flushleft} \large
            Alunos:\\
            Guilherme S. L. Lopes -- RA143630\\
            Lucas Nakamura Batista -- RA143494\\
        \end{flushleft}
    \end{minipage}
    ~
    \begin{minipage}{0.4\textwidth}
        \begin{flushright} \large
            Professor: \\
            Ronaldo A. L. Gonçalves
        \end{flushright}
    \end{minipage}
    
    \vfill % Pushes the remaining content to the bottom of the page
    
    % 5. Date / Footer
    \large \today \\
\end{titlepage}

\section{Funcionamento Geral e Modular}
\subsection{Funcionamento Geral}
O trabalho realizado consiste em uma "calculadora" de determinantes de matrizes 
de qualquer ordem. Nosso programa funciona utilizando de uma leitura de um 
arquivo .txt com uma tabela de números, e ao selecionar uma ordem $m$, pegamos 
as $m$ linhas e $m$ colunas dessa tabela, e fazemos o cálculo do determinante 
da matriz resultante dessa seleção.

Após realizar a seleção da matriz, o calculo do determinante é feito de forma 
simples, usando do Teorema de Laplace e da Regra de Sarrus. Ao jogar a matriz 
no software, ele compara a ordem, caso ela seja de um valor superior a quatro, 
será feito o cálculo usando o Teorema de Laplace, e caso seja menor, fará ou 
por Regra de Sarrus, ou caso for de ordem 2 ou 1, fará o calculo do determinante 
usando da multiplicação da diagonal principal menos a diagonal secundária, ou 
retornará o numéro da matriz unidade.

Caso a matriz estiver na condição para calcular utilizando do Laplace, ela cria 
uma matriz resultande da exclusão da linha e da coluna da matriz original, e 
calcula o determinante dessa matriz determinante de forma recursiva, até chegar 
em uma matriz de ordem 3. Após calcular o determinante da matriz resultante, é 
feito o calculo do cofator, multiplicando o determinante por $-1$ elevado a 
linha mais a coluna do elemento do cofator.

\subsection{Funcionamento Modular}
\begin{itemize}
    \item \verb|Matriz criaMatriz(int ordem);|
    \begin{itemize}
        \item Recebe como parâmetro um inteiro representando a ordem da matriz, 
        aloca dinamicamente na memória os espaços para a linha e coluna do 
        tamanho da \verb|ordem| e retorna essa matriz.
    \end{itemize}
    \item \verb|void alocaValoresMatriz(Matriz *matriz, int ordem);|
    \begin{itemize}
        \item Recebe como parâmetro uma matriz vazia e a ordem dessa matriz, lê 
        o arquivo "\verb|matriz.txt|" que contém diversos valores diferentes 
        que estão ordenados por linha e coluna, verifica a ordem da matriz e 
        aloca os elementos de suas respectivas posições.
    \end{itemize}
    \item \verb|double calculaMatriz(Matriz m, int ordem);|
    \begin{itemize}
        \item Recebe uma matriz e calcula seu determinante. Se a ordem for maior 
        que 3, ela chama o cálculo por Teorema de Laplace. Se a ordem for 3, 
        usa a Regra de Sarrus. Se for 2, aplica a fórmula direta $a_{11}a_{22} - 
        a_{12}a_{21}$. Se for 1, retorna o único elemento da matriz.
    \end{itemize}
    \item \verb|double calculaSarrus(Matriz m);|
    \begin{itemize}
        \item Define dois valores, um com a soma dos produtos das diagonais 
        principais e outro com a soma dos produtos das diagonais secundárias.
        \[
        \left\{
        \begin{array}{l}
        a = (a_{11} \cdot a_{22} \cdot a_{33}) + (a_{12} \cdot a_{23} \cdot 
        a_{31}) + (a_{13} \cdot a_{21} \cdot a_{32}) \\
        b = (a_{31} \cdot a_{22} \cdot a_{13}) + (a_{32} \cdot a_{23} \cdot 
        a_{11}) + (a_{33} \cdot a_{21} \cdot a_{12})
        \end{array}
        \right.
        \]
        Depois retorna o valor da subtração de $a$ por $b$ ($a-b$)
    \end{itemize}
    \item \verb|double calculaLaplace(Matriz mat, int ordem);|
    \begin{itemize}
        \item Cria uma matriz resultante da exclusão da linha $i$ e da coluna 
        $j$, após isso calcula o determinante dessa matriz resultante, multiplica 
        pelo item $a_{ij}$ da matriz original, e multiplica por $(-1)^{i + j}$.
    \end{itemize}
    \item \verb|Matriz matrizResultante(Matriz mat, int i, int j, int ordem);|
    \begin{itemize}
        \item Cria a matriz resultante excluindo a linha $i$ e a coluna $j$, mas 
        mantendo os outros elementos.
    \end{itemize}
    \item \verb|void exibirMatriz(Matriz matriz, int ordem);|
    \begin{itemize}
        \item Exibe a matriz selecionada pelo usuário formatada, com os números 
        em suas respectivas linhas e colunas.
    \end{itemize}
    \item \verb|void esvaziaMemoria(Matriz matriz, int ordem);|
    \begin{itemize}
        \item Esvazia a memória de uma matriz, percorrendo os vetores iniciais 
        dando free neles, e no fim dando um free na matriz.
    \end{itemize}
\end{itemize}

\section{Atividades dos Participantes}
No trabalho, as contribuições dos alunos Guilherme e Lucas foram as implementações 
das seguintes funções: 
\subsection{Guilherme dos Santos Lira Lopes}
\begin{itemize}
    \item \verb|alocaValoresMatriz|
    \item \verb|calculaMatriz|  
    \item \verb|exibirMatriz|
    \item \verb|calculaSarrus|
    \item \verb|esvaziaMemoria|
\end{itemize}

\subsection{Lucas Nakamura Batista}
\begin{itemize}
    \item \verb|criaMatriz|
    \item \verb|calculaLaplace|
    \item \verb|matrizResultante|
    \item \verb|potencia|
\end{itemize}

Ao fim da implementação das funções, nos juntamos para corrigir quaisquer erros 
restantes, por exemplo o código não funcionava corretamete depois de uma certa 
ordem, ou ele retornava um valor próximo do resultado correto devido a diferença 
entre o float -- o que estávamos usando até então -- e o double.

\section{Comentários Gerais}
Nesse trabalho, nossa maior dificuldade foi na implementação do calculo de Laplace, 
visto que é um processo trabalhoso, e também a correção de alguns erros, como 
quando uma matriz era de uma ordem acima de 5, essa matriz não era alocada 
corretamente, e também o supramencionado erro de arredondamento do float, que 
deixava o resultado incorreto. 

De certa forma não achamos que o trabalho tinha uma dificuldade tão elevada, 
com exceção de Laplace e na correção de erros, além disso devido a nossa 
implementação, o tempo de execução para calcular matrizes de ordem acima de 11 
fica extremamente elevado e não é possivel obter o resultado, essa é a única 
coisa que encontramos que não funciona dentro do nosso código.

\end{document}